% Library unit: the short side (rook endings). Printed once, before the first position that uses it; see lib/README.md.
\paragraph{Short side}\label{short-side}\index{Rook!Short side}

The narrower part is the
\term{short-side}.
In \xref{diagram}{side-checks} the e-pawn has three files on its right (f, g, h).

The defending king belongs on the short side. There it stays near the pawn without standing in
the way of its own rook, which gives side checks from the \term{long-side}.
